
\section{Contexte}
\ifprof
\else
Dans la volonté de rééducation à la marche des patients partiellement ou totalement paralysés au niveau des membres inférieurs, les exosquelettes de marche offrent une solution de plus en plus envisagée par les structures de soins. Si les technologies existantes permettent des séances de rééducation musculaire efficaces, elles n'offrent pour le moment que peu de possibilités d'amélioration dans la vie quotidienne des patients, du fait de la nécessité de la présence d'un personnel soignant formé à leur utilisation.\\
Il apparait ainsi qu'une solution permettant une plus grande autonomie des patients, que ce soit pour des séances de rééducation encadrées, ou pour leurs déplacements dans la vie quotidienne serait une grande avancée. Elle permettrait ainsi des fréquences, intensités et durées de séances plus en phase avec la pathologie de la personne équipée, tout en facilitant sa mobilité au quotidien.\\
Pensé pour minimiser le temps de formation pour le patient et le thérapeute, l'exosquelette Atalante (figure \ref{2023_ccs_psi_fig_01}), développé par la société Wandercraft a pour volonté d'optimiser les séances de rééducation, et à terme d'aboutir à une solution permettant l'autonomie quasi totale du patient. En effet, ce système permet la verticalisation et des déplacements pouvant s'affranchir de toute dépendance à une tierce personne. De par sa liberté d'utilisation pour le patient, les bénéfices sont importants : possibilité de retours sensoriels, flexibilité de l'entrainement à la marche et à la course ou encore personnalisation des programmes proposés.\\

\begin{center}
\includegraphics[width=.6\textwidth]{2025_07_03_97545f5dc188959e5663g-01}
\captionof{figure}{Exosquelette Atalante et modélisation 3D associée \label{2023_ccs_psi_fig_01}}
%Figure 1 
\end{center}

L'architecture de l'exosquelette autorise la complète autonomie des membres supérieurs du patient. De plus, le réglage rapide des dimensions des jambes et des tibias de l'exosquelette et une utilisation intuitive de son paramétrage permettent une utilisation plus aisée. Ainsi, associé avec la possibilité de variation de l'aide apportée par les différents actionneurs, cet équipement est en parfaite adéquation avec l'enchainement des séances de rééducation où les déplacements sont totalement pris en charge, quelles que soient la pathologie et la morphologie du patient.\\
Les exercices programmés par les thérapeutes s'articulent autour de deux stratégies thérapeutiques différentes, selon le handicap du patient :

\begin{itemize}
  \item la proprioception (sa propre perception) de la verticalité et de la marche pour les patients paraplégiques ou ayant des problèmes d'équilibre ;
  \item le renforcement musculaire pour les patients ayant subi un traumatisme important.
\end{itemize}

Afin de permettre au patient de reproduire une marche comparable à celle d'un humain valide, et ce quel que soit l'exercice de rééducation préprogrammé, l'exosquelette Atalante possède douze degrés de liberté, six par jambe, et chaque mobilité est contrôlée par un actionneur électrique.\\
L'ensemble de l'étude menée dans ce sujet se limitera à une marche en ligne droite (figure \ref{2023_ccs_psi_fig_03}), dans le plan sagittal tel que défini figure \ref{2023_ccs_psi_fig_02}.\\
La loi de commande des actionneurs associés aux différentes articulations est donc conditionnée par :

\begin{itemize}
  \item les trajectoires des différents membres lors de la marche, prédéterminées par un serveur extérieur lors de la prise de mesures du patient ;
  \item la stabilisation du patient lors de la marche ;
  \item la rééducation musculaire du patient, selon un protocole dicté par le thérapeute.\\
\end{itemize}
Le cahier des charges partiel de l'exosquelette à vérifier dans ce sujet est fourni en annexe A.\\


\begin{center}
\includegraphics[width=.4\textwidth]{2025_07_03_97545f5dc188959e5663g-02}
\captionof{figure}{Description des plans d'évolution \label{2023_ccs_psi_fig_02}}
%Figure 2 
\end{center}


\begin{center}
\includegraphics[width=\textwidth]{2025_07_03_97545f5dc188959e5663g-02(1)}
\captionof{figure}{Description d'une foulée en ligne droite \label{2023_ccs_psi_fig_03}}
%Figure 3 
\end{center}

\fi